% Options for packages loaded elsewhere
\PassOptionsToPackage{unicode}{hyperref}
\PassOptionsToPackage{hyphens}{url}
%
\documentclass[
  10pt,
  ignorenonframetext,
]{beamer}
\usepackage{pgfpages}
\setbeamertemplate{caption}[numbered]
\setbeamertemplate{caption label separator}{: }
\setbeamercolor{caption name}{fg=normal text.fg}
\beamertemplatenavigationsymbolsempty
% Prevent slide breaks in the middle of a paragraph
\widowpenalties 1 10000
\raggedbottom
\setbeamertemplate{part page}{
  \centering
  \begin{beamercolorbox}[sep=16pt,center]{part title}
    \usebeamerfont{part title}\insertpart\par
  \end{beamercolorbox}
}
\setbeamertemplate{section page}{
  \centering
  \begin{beamercolorbox}[sep=12pt,center]{part title}
    \usebeamerfont{section title}\insertsection\par
  \end{beamercolorbox}
}
\setbeamertemplate{subsection page}{
  \centering
  \begin{beamercolorbox}[sep=8pt,center]{part title}
    \usebeamerfont{subsection title}\insertsubsection\par
  \end{beamercolorbox}
}
\AtBeginPart{
  \frame{\partpage}
}
\AtBeginSection{
  \ifbibliography
  \else
    \frame{\sectionpage}
  \fi
}
\AtBeginSubsection{
  \frame{\subsectionpage}
}
\usepackage{amsmath,amssymb}
\usepackage{iftex}
\ifPDFTeX
  \usepackage[T1]{fontenc}
  \usepackage[utf8]{inputenc}
  \usepackage{textcomp} % provide euro and other symbols
\else % if luatex or xetex
  \usepackage{unicode-math} % this also loads fontspec
  \defaultfontfeatures{Scale=MatchLowercase}
  \defaultfontfeatures[\rmfamily]{Ligatures=TeX,Scale=1}
\fi
\usepackage{lmodern}
\usetheme[]{SimplePlus}
\ifPDFTeX\else
  % xetex/luatex font selection
\fi
% Use upquote if available, for straight quotes in verbatim environments
\IfFileExists{upquote.sty}{\usepackage{upquote}}{}
\IfFileExists{microtype.sty}{% use microtype if available
  \usepackage[]{microtype}
  \UseMicrotypeSet[protrusion]{basicmath} % disable protrusion for tt fonts
}{}
\makeatletter
\@ifundefined{KOMAClassName}{% if non-KOMA class
  \IfFileExists{parskip.sty}{%
    \usepackage{parskip}
  }{% else
    \setlength{\parindent}{0pt}
    \setlength{\parskip}{6pt plus 2pt minus 1pt}}
}{% if KOMA class
  \KOMAoptions{parskip=half}}
\makeatother
\usepackage{xcolor}
\newif\ifbibliography
\usepackage{graphicx}
\makeatletter
\def\maxwidth{\ifdim\Gin@nat@width>\linewidth\linewidth\else\Gin@nat@width\fi}
\def\maxheight{\ifdim\Gin@nat@height>\textheight\textheight\else\Gin@nat@height\fi}
\makeatother
% Scale images if necessary, so that they will not overflow the page
% margins by default, and it is still possible to overwrite the defaults
% using explicit options in \includegraphics[width, height, ...]{}
\setkeys{Gin}{width=\maxwidth,height=\maxheight,keepaspectratio}
% Set default figure placement to htbp
\makeatletter
\def\fps@figure{htbp}
\makeatother
\setlength{\emergencystretch}{3em} % prevent overfull lines
\providecommand{\tightlist}{%
  \setlength{\itemsep}{0pt}\setlength{\parskip}{0pt}}
\setcounter{secnumdepth}{-\maxdimen} % remove section numbering
% From https://tex.stackexchange.com/questions/75667/change-colour-on-chapter-section-headings-lyx
\usepackage{xcolor}
\definecolor{darkgreen}{rgb}{0,0.5,0}

\usepackage{cancel}

%\usepackage{sectsty}
%\chapterfont{\color{blue}}
%\sectionfont{\color{red}}
%\subsectionfont{\color{purple}}
%\subsubsectionfont{\color{blue}}

% From https://tex.stackexchange.com/questions/1455/how-to-set-the-font-for-a-section-title-and-chapter-etc
%\usepackage{titlesec}
%\titleformat{\chapter}[display]
%  {\normalfont\sffamily\huge\bfseries\color{blue}}
%  {\chaptertitlename\ \thechapter}{20pt}{\Huge}
%\titleformat{\section}
%  {\normalfont\sffamily\Large\bfseries\color{cyan}}
%  {\thesection}{1em}{}

% Work-around for too deeply nested error.
\usepackage{enumitem}
\setlistdepth{9}
\setlist[itemize,1]{label=\textendash}
\setlist[itemize,2]{label=\textendash}
\setlist[itemize,3]{label=\textendash}
\setlist[itemize,4]{label=\textendash}
\setlist[itemize,5]{label=\textendash}
\setlist[itemize,6]{label=\textendash}
\setlist[itemize,7]{label=\textendash}
\setlist[itemize,8]{label=\textendash}
\setlist[itemize,9]{label=\textendash}
\renewlist{itemize}{itemize}{9}

% \setlist[itemize,1]{label=\textbullet}
% \setlist[itemize,2]{label=\textendash}
% \setlist[itemize,3]{label=\textasteriskcentered}
% \setlist[itemize,4]{label=\textperiodcentered}
% setlist[itemize,5]{label=\textbullet}
% \setlist[itemize,6]{label=\textbullet}
% \setlist[itemize,7]{label=\textbullet}
% \setlist[itemize,8]{label=\textbullet}
% \setlist[itemize,9]{label=\textbullet}

% Configure enumerate lists to prevent infinite recursion in Beamer themes.
\setlist[enumerate,1]{label=\arabic*.,ref=\arabic*}
\setlist[enumerate,2]{label=\alph*.,ref=\theenumi.\alph*}
\setlist[enumerate,3]{label=\roman*.,ref=\theenumi.\theenumii.\roman*}
\setlist[enumerate,4]{label=\Alph*.,ref=\theenumi.\theenumii.\theenumiii.\Alph*}
\setlist[enumerate,5]{label=\Roman*.,ref=\theenumi.\theenumii.\theenumiii.\theenumiv.\Roman*}
\setlist[enumerate,6]{label=\arabic*.,ref=\theenumi.\theenumii.\theenumiii.\theenumiv.\theenumv.\arabic*}
\setlist[enumerate,7]{label=\alph*.,ref=\theenumi.\theenumii.\theenumiii.\theenumiv.\theenumv.\theenumvi.\alph*}
\setlist[enumerate,8]{label=\roman*.,ref=\theenumi.\theenumii.\theenumiii.\theenumiv.\theenumv.\theenumvi.\theenumvii.\roman*}
\setlist[enumerate,9]{label=\arabic*.,ref=\theenumi.\theenumii.\theenumiii.\theenumiv.\theenumv.\theenumvi.\theenumvii.\theenumviii.\arabic*}
\renewlist{enumerate}{enumerate}{9}

% #############################################################################
% Vector / matrix notation
% #############################################################################

% \vec is already defined so we use \vv

% Use bold.
%\newcommand{\vv}[1]{\boldsymbol{#1}}
%\newcommand{\mat}[1]{\boldsymbol{#1}}

% Use underline.
% From https://latex.org/forum/viewtopic.php?t=17798
% ulem doesn't work well in math mode.
%\newcommand{\vv}[1]{\boldsymbol{\uline{#1}}}
%\newcommand{\mat}[1]{\boldsymbol{\uuline{#1}}}
\newcommand{\vv}[1]{\boldsymbol{\underline{#1}}}
\newcommand{\mat}[1]{\boldsymbol{\underline{\underline{#1}}}}

% TODO(saggese): Remove this and use \quad.
\newcommand{\eqspace}{\hspace{10 mm}}
\newcommand{\norm}[1]{\left\lVert#1\right\rVert}

% NOTE: These shortcuts are not only to make the latex code easier but also to
% make easy to find out what the standard representation is.
% In general we prefer to put \vv{} and \mat{} at outermost.

% #############################################################################
% Newshortcuts for matrices and vectors.
% #############################################################################

% - v* for vectors, e.g., va, vb, ..., vmu, vA, vGamma
% - m* for matrices, e.g., mA, ..., mSigma

% Vector.
\newcommand{\vA}{\vv{A}}
\newcommand{\vB}{\vv{B}}
\newcommand{\vC}{\vv{C}}
\newcommand{\vD}{\vv{D}}
\newcommand{\vDelta}{\vv{\Delta}}
\newcommand{\vE}{\vv{E}}
\newcommand{\vF}{\vv{F}}
\newcommand{\vG}{\vv{G}}
\newcommand{\vGamma}{\vv{\Gamma}}
\newcommand{\vH}{\vv{H}}
\newcommand{\vI}{\vv{I}}
\newcommand{\vJ}{\vv{J}}
\newcommand{\vK}{\vv{K}}
\newcommand{\vL}{\vv{L}}
\newcommand{\vLambda}{\vv{\Lambda}}
\newcommand{\vM}{\vv{M}}
\newcommand{\vN}{\vv{N}}
\newcommand{\vO}{\vv{O}}
\newcommand{\vOmega}{\vv{\Omega}}
\newcommand{\vP}{\vv{P}}
\newcommand{\vPhi}{\vv{\Phi}}
\newcommand{\vPi}{\vv{\Pi}}
\newcommand{\vPsi}{\vv{\Psi}}
\newcommand{\vQ}{\vv{Q}}
\newcommand{\vR}{\vv{R}}
\newcommand{\vS}{\vv{S}}
\newcommand{\vSigma}{\vv{\Sigma}}
\newcommand{\vT}{\vv{T}}
\newcommand{\vU}{\vv{U}}
\newcommand{\vUpsilon}{\vv{\Upsilon}}
\newcommand{\vV}{\vv{V}}
\newcommand{\vW}{\vv{W}}
\newcommand{\vX}{\vv{X}}
\newcommand{\vXi}{\vv{\Xi}}
\newcommand{\vY}{\vv{Y}}
\newcommand{\vZ}{\vv{Z}}
\newcommand{\va}{\vv{a}}
\newcommand{\valpha}{\vv{\alpha}}
\newcommand{\vb}{\vv{b}}
\newcommand{\vbeta}{\vv{\beta}}
\newcommand{\vc}{\vv{c}}
\newcommand{\vchi}{\vv{\chi}}
\newcommand{\vd}{\vv{d}}
\newcommand{\vdelta}{\vv{\delta}}
\newcommand{\ve}{\vv{e}}
\newcommand{\vepsilon}{\vv{\epsilon}}
\newcommand{\veta}{\vv{\eta}}
\newcommand{\vf}{\vv{f}}
\newcommand{\vg}{\vv{g}}
\newcommand{\vgamma}{\vv{\gamma}}
\newcommand{\vh}{\vv{h}}
\newcommand{\vi}{\vv{i}}
\newcommand{\viota}{\vv{\iota}}
\newcommand{\vj}{\vv{j}}
\newcommand{\vk}{\vv{k}}
\newcommand{\vkappa}{\vv{\kappa}}
\newcommand{\vl}{\vv{l}}
\newcommand{\vlambda}{\vv{\lambda}}
\newcommand{\vm}{\vv{m}}
\newcommand{\vmu}{\vv{\mu}}
\newcommand{\vn}{\vv{n}}
\newcommand{\vnu}{\vv{\nu}}
\newcommand{\vo}{\vv{o}}
\newcommand{\vomega}{\vv{\omega}}
\newcommand{\vp}{\vv{p}}
\newcommand{\vphi}{\vv{\phi}}
\newcommand{\vpi}{\vv{\pi}}
\newcommand{\vpsi}{\vv{\psi}}
\newcommand{\vq}{\vv{q}}
\newcommand{\vr}{\vv{r}}
\newcommand{\vrho}{\vv{\rho}}
\newcommand{\vs}{\vv{s}}
\newcommand{\vsigma}{\vv{\sigma}}
\newcommand{\vt}{\vv{t}}
\newcommand{\vtau}{\vv{\tau}}
\newcommand{\vtheta}{\vv{\theta}}
\newcommand{\vu}{\vv{u}}
\newcommand{\vupsilon}{\vv{\upsilon}}
\newcommand{\vvarepsilon}{\vv{\varepsilon}}
\newcommand{\vvarphi}{\vv{\varphi}}
\newcommand{\vvarrho}{\vv{\varrho}}
\newcommand{\vvartheta}{\vv{\vartheta}}
\newcommand{\vvv}{\vv{v}}
\newcommand{\vw}{\vv{w}}
\newcommand{\vx}{\vv{x}}
\newcommand{\vxi}{\vv{\xi}}
\newcommand{\vy}{\vv{y}}
\newcommand{\vz}{\vv{z}}
\newcommand{\vzeta}{\vv{\zeta}}

% Matrix.
% \mat -> \mm?
\newcommand{\mA}{\mat{A}}
\newcommand{\mB}{\mat{B}}
\newcommand{\mC}{\mat{C}}
\newcommand{\mD}{\mat{D}}
\newcommand{\mDelta}{\mat{\Delta}}
\newcommand{\mE}{\mat{E}}
\newcommand{\mF}{\mat{F}}
\newcommand{\mG}{\mat{G}}
\newcommand{\mGamma}{\mat{\Gamma}}
\newcommand{\mH}{\mat{H}}
\newcommand{\mI}{\mat{I}}
\newcommand{\mJ}{\mat{J}}
\newcommand{\mK}{\mat{K}}
\newcommand{\mL}{\mat{L}}
\newcommand{\mLambda}{\mat{\Lambda}}
\newcommand{\mM}{\mat{M}}
\newcommand{\mN}{\mat{N}}
\newcommand{\mO}{\mat{O}}
\newcommand{\mOmega}{\mat{\Omega}}
\newcommand{\mP}{\mat{P}}
\newcommand{\mPi}{\mat{\Pi}}
\newcommand{\mPsi}{\mat{\Psi}}
\newcommand{\mQ}{\mat{Q}}
\newcommand{\mR}{\mat{R}}
\newcommand{\mS}{\mat{S}}
\newcommand{\mSigma}{\mat{\Sigma}}
\newcommand{\mT}{\mat{T}}
\newcommand{\mU}{\mat{U}}
\newcommand{\mUpsilon}{\mat{\Upsilon}}
\newcommand{\mV}{\mat{V}}
\newcommand{\mW}{\mat{W}}
\newcommand{\mX}{\mat{X}}
\newcommand{\mXi}{\mat{\Xi}}
\newcommand{\mY}{\mat{Y}}
\newcommand{\mZ}{\mat{Z}}

% #############################################################################
% Other shortcuts
% #############################################################################

% Number fields
% TODO: use \R, \N, ... and the standard symbols (see
% http://oeis.org/wiki/List_of_LaTeX_mathematical_symbols)
\newcommand{\bbB}{\mathbb{B}}
\newcommand{\bbC}{\mathbb{C}}
\newcommand{\bbF}{\mathbb{F}}
\newcommand{\bbL}{\mathbb{L}}
\newcommand{\bbN}{\mathbb{N}}
\newcommand{\bbQ}{\mathbb{Q}}
\newcommand{\bbR}{\mathbb{R}}
\newcommand{\bbS}{\mathbb{S}}
\newcommand{\bbZ}{\mathbb{Z}}

% Script vars.
\newcommand{\cala}{\mathcal{a}}
\newcommand{\calb}{\mathcal{b}}
\newcommand{\calc}{\mathcal{c}}
\newcommand{\cald}{\mathcal{d}}
\newcommand{\cale}{\mathcal{e}}
\newcommand{\calf}{\mathcal{f}}
\newcommand{\calg}{\mathcal{g}}
\newcommand{\calh}{\mathcal{h}}
\newcommand{\cali}{\mathcal{i}}
\newcommand{\calj}{\mathcal{j}}
\newcommand{\calk}{\mathcal{k}}
\newcommand{\call}{\mathcal{l}}
\newcommand{\calm}{\mathcal{m}}
\newcommand{\caln}{\mathcal{n}}
\newcommand{\calo}{\mathcal{o}}
\newcommand{\calp}{\mathcal{p}}
\newcommand{\calq}{\mathcal{q}}
\newcommand{\calr}{\mathcal{r}}
\newcommand{\cals}{\mathcal{s}}
\newcommand{\calt}{\mathcal{t}}
\newcommand{\calu}{\mathcal{u}}
\newcommand{\calv}{\mathcal{v}}
\newcommand{\calw}{\mathcal{w}}
\newcommand{\calx}{\mathcal{x}}
\newcommand{\caly}{\mathcal{y}}
\newcommand{\calz}{\mathcal{z}}
\newcommand{\calA}{\mathcal{A}}
\newcommand{\calB}{\mathcal{B}}
\newcommand{\calC}{\mathcal{C}}
\newcommand{\calD}{\mathcal{D}}
\newcommand{\calE}{\mathcal{E}}
\newcommand{\calF}{\mathcal{F}}
\newcommand{\calG}{\mathcal{G}}
\newcommand{\calH}{\mathcal{H}}
\newcommand{\calI}{\mathcal{I}}
\newcommand{\calJ}{\mathcal{J}}
\newcommand{\calK}{\mathcal{K}}
\newcommand{\calL}{\mathcal{L}}
\newcommand{\calM}{\mathcal{M}}
\newcommand{\calN}{\mathcal{N}}
\newcommand{\calO}{\mathcal{O}}
\newcommand{\calP}{\mathcal{P}}
\newcommand{\calQ}{\mathcal{Q}}
\newcommand{\calR}{\mathcal{R}}
\newcommand{\calS}{\mathcal{S}}
\newcommand{\calT}{\mathcal{T}}
\newcommand{\calU}{\mathcal{U}}
\newcommand{\calV}{\mathcal{V}}
\newcommand{\calW}{\mathcal{W}}
\newcommand{\calX}{\mathcal{X}}
\newcommand{\calY}{\mathcal{Y}}
\newcommand{\calZ}{\mathcal{Z}}

% tilde
% TODO(gp): Expand this.
\newcommand{\ttau}{\tilde{\tau}}
% Irregular to avoid conflict.
\newcommand{\ttt}{\tilde{t}}
\newcommand{\ty}{\tilde{y}}

% TODO(gp): Expand this.
\newcommand{\hF}{\widehat{F}}
\newcommand{\hvF}{\vv{\widehat{F}}}
\newcommand{\hG}{\widehat{G}}
\newcommand{\hvG}{\vv{\widehat{G}}}
\newcommand{\hX}{\widehat{X}}
\newcommand{\hvX}{\vv{\widehat{X}}}

% Unit vector.
% Same as the vector name with unit
% TODO(gp): Merge with above.
\newcommand{\vvvhat}{\vv{\hat{v}}}
\newcommand{\uuhat}{\vv{\hat{u}}}
\newcommand{\xxhat}{\vv{\hat{x}}}
%
\newcommand{\tildeAAA}{\mat{\tilde{A}}}

% Basic.

% Define.
%\def\defeq{\overset{\begin{tiny}def\end{tiny}}{=}}
\def\defeq{\triangleq}
% TODO: This is already defined, but not clear where
%\newcommand{\gcd}{\text{gcd}}
\newcommand{\lcm}{\text{lcm}}

% Mean and variance.
\newcommand{\EE}{\mathbb{E}}
\newcommand{\VV}{\mathbb{V}}
\newcommand{\Cov}{\text{Cov}}
\newcommand{\Corr}{\text{Corr}}
%\newcommand{\Pr}{\mathbb{P}}

% Distributions.
% text{} or textsf
\def\Bernoulli{\textsf{Bernoulli}}
\def\Beta{\textsf{Beta}}
\def\Binomial{\textsf{Binomial}}
\def\Exponential{\textsf{Exponential}}
\def\LogNormal{\textsf{LogNormal}}
\def\N{\textsf{N}}
\def\Poisson{\textsf{Poisson}}
\def\Uniform{\textsf{Uniform}}
\def\WN{\textsf{WN}}
\def\GWN{\textsf{GWN}}

% Linear algebra.
\def\Col{\textsf{Col}}
\def\Dim{\textsf{Dim}}
\def\Re{\textsf{Re}}
\def\Im{\textsf{Im}}
\def\Ker{\textsf{Ker}}
\def\Null{\textsf{Null}}
\def\Rank{\textsf{Rank}}
\def\Row{\textsf{Row}}
\def\Span{\textsf{Span}}
\def\Det{\textsf{Det}}
\def\Tr{\textsf{Tr}}
\def\Diag{\textsf{Diag}}
\def\sign{\textsf{sign}}
\def\argmin{\textsf{argmin}}
\def\argmax{\textsf{argmax}}
\def\mod{\textsf{ mod }}
\def\min{\textsf{min}}
\def\max{\textsf{max}}
\def\avg{\textsf{avg}}

% Derivative at a point
% https://tex.stackexchange.com/questions/64176/derivative-at-a-point
\newcommand{\at}[2][]{#1|_{#2}}

% Logic operators.
\def\NOT{\textsf{\begin{footnotesize}NOT\end{footnotesize}}}
\def\OR{\textsf{\begin{footnotesize} OR \end{footnotesize}}}
\def\AND{\textsf{\begin{footnotesize} AND \end{footnotesize}}}
\def\XOR{\textsf{\begin{footnotesize} XOR \end{footnotesize}}}
\def\IFF{\textsf{\begin{footnotesize} IFF \end{footnotesize}}}

% For some reasons new commands can't end with a number.
%\newcommand{\v0}{\vv{0}}
%\newcommand{\v1}{\vv{1}}
\ifLuaTeX
  \usepackage{selnolig}  % disable illegal ligatures
\fi
\usepackage{bookmark}
\IfFileExists{xurl.sty}{\usepackage{xurl}}{} % add URL line breaks if available
\urlstyle{same}
\hypersetup{
  hidelinks,
  pdfcreator={LaTeX via pandoc}}

\author{}
\date{}

\begin{document}

\begin{frame}
\let\emph\textit
\let\uline\underline
\let\ul\underline

\begin{columns}[T]
\begin{column}{0.15\textwidth}
\includegraphics{msml610/lectures_source/figures/UMD_Logo.png}
\end{column}

\begin{column}{0.75\textwidth}
\vspace{0.4cm}
\begingroup \large

MSML610: Advanced Machine Learning \endgroup
\end{column}
\end{columns}

\vspace{1cm}

\begingroup \Large

\textbf{\[\text{\textcolor{blue}{Lesson 11.1: Decision-Making with Causal Models}}\]}
\endgroup

\vspace{1cm}

\begin{columns}[T]
\begin{column}{0.65\textwidth}
\textbf{Instructor}: Dr.~GP Saggese,
\href{gsaggese@umd.edu}{\textcolor{blue}{\underline{gsaggese@umd.edu}}}

\textbf{References}:
\end{column}

\begin{column}{0.4\textwidth}
\end{column}
\end{columns}
\end{frame}

\begin{frame}{}
\phantomsection\label{section}
\begin{itemize}
\tightlist
\item
  \emph{\textbf{\textcolor{red}{Why Prediction Is Not Enough}}}

  \begin{itemize}
  \tightlist
  \item
    Prediction Pipelines vs.~Decision Pipelines
  \item
    When Prediction Fails
  \item
    Causal Models as the Foundation for Decisions
  \end{itemize}
\item
  Foundations of Decision Theory
\item
  Decision Support with Causal Models
\item
  Sequential Decision-Making and Active Learning
\item
  Uncertainty in Causal Decisions
\end{itemize}
\end{frame}

\begin{frame}{}
\phantomsection\label{section-1}
\begin{itemize}
\tightlist
\item
  Why Prediction Is Not Enough

  \begin{itemize}
  \tightlist
  \item
    \emph{\textbf{\textcolor{red}{Prediction Pipelines vs. Decision Pipelines}}}
  \item
    When Prediction Fails
  \item
    Causal Models as the Foundation for Decisions
  \end{itemize}
\item
  Foundations of Decision Theory
\item
  Decision Support with Causal Models
\item
  Sequential Decision-Making and Active Learning
\item
  Uncertainty in Causal Decisions
\end{itemize}
\end{frame}

\begin{frame}{From Predictions to Decisions}
\phantomsection\label{from-predictions-to-decisions}
\begin{itemize}
\tightlist
\item
  \textbf{\textcolor{red}{Prediction}} answers: \emph{``What will
  happen?''}

  \begin{itemize}
  \tightlist
  \item
    Input: features \(X\)
  \item
    Output: estimate \(\hat{y} = \EE[Y | X]\)
  \item
    Optimized for accuracy on a fixed distribution
  \end{itemize}
\item
  \textbf{\textcolor{orange}{Decision}} answers: \emph{``What should I
  do?''}

  \begin{itemize}
  \tightlist
  \item
    Input: current state, set of possible actions \(\mathcal{A}\)
  \item
    Output: action \(a^* \in \mathcal{A}\) that maximizes expected
    utility
  \item
    Requires reasoning about
    \textbf{\textcolor{darkgreen}{what would change}} under each action
  \end{itemize}
\item
  \textbf{\textcolor{teal}{Q}}: If we already have a highly accurate
  predictor, isn't the decision obvious?

  \begin{itemize}
  \tightlist
  \item
    \textbf{\textcolor{cyan}{No}} --- a good predictor tells us what
    correlates with good outcomes, not what \emph{causes} them
  \end{itemize}
\end{itemize}
\end{frame}

\begin{frame}{Prediction Pipeline: Standard Supervised Learning}
\phantomsection\label{prediction-pipeline-standard-supervised-learning}
\begin{itemize}
\tightlist
\item
  Typical ML system is optimized \emph{end-to-end} for prediction:

  \begin{enumerate}
  \tightlist
  \item
    Collect observational data \((X_i, Y_i)\)
  \item
    Fit a model \(\hat{f}(X) \approx \EE[Y | X]\)
  \item
    Deploy to score new units
  \end{enumerate}
\item
  \textbf{\textcolor{red}{Crucial assumption}}: the joint distribution
  \(\Pr(X, Y)\) at deployment is the same as at training

  \begin{itemize}
  \tightlist
  \item
    No intervention changes how \(X\) is generated
  \item
    No agent \emph{acts} on the prediction
  \end{itemize}
\item
  \textbf{\textcolor{blue}{Example}}: a spam filter trained on past
  emails to predict \(\Pr(\text{spam} | X)\)

  \begin{itemize}
  \tightlist
  \item
    Accurate predictor \(\to\) useful as long as spammers do not adapt
  \end{itemize}
\end{itemize}
\end{frame}

\begin{frame}{Decision Pipeline: What Changes}
\phantomsection\label{decision-pipeline-what-changes}
\begin{itemize}
\tightlist
\item
  A \textbf{\textcolor{red}{decision pipeline}} explicitly represents:

  \begin{itemize}
  \tightlist
  \item
    The set of available actions \(\mathcal{A}\)
  \item
    A utility function \(U(y, a)\) over outcomes and actions
  \item
    A \textbf{\textcolor{orange}{causal}} model
    \(\Pr(Y | \operatorname{do}(a))\) of how actions affect outcomes
  \end{itemize}
\item
  \textbf{\textcolor{darkgreen}{Key distinction}}: \(\Pr(Y | X = x)\) is
  observational; \(\Pr(Y | \operatorname{do}(X = x))\) is interventional

  \begin{itemize}
  \tightlist
  \item
    These can differ drastically when confounders are present
  \item
    Decisions depend on the \textbf{\textcolor{teal}{interventional}}
    quantity
  \end{itemize}
\item
  \textbf{\textcolor{cyan}{Example}}: a hospital wants to decide
  \emph{whether to prescribe drug A}

  \begin{itemize}
  \tightlist
  \item
    Observational: \(\Pr(\text{recovery} | \text{took drug A})\)
  \item
    Interventional:
    \(\Pr(\text{recovery} | \operatorname{do}(\text{give drug A}))\)
  \item
    The first is biased by who chose to take the drug
  \end{itemize}
\end{itemize}
\end{frame}

\begin{frame}{Prediction vs.~Decision: Side-by-Side}
\phantomsection\label{prediction-vs.-decision-side-by-side}
\begin{columns}[T]
\begin{column}{0.5\textwidth}
\textbf{\textcolor{red}{Prediction}}

\begin{itemize}
\tightlist
\item
  Target: \(\EE[Y | X]\)
\item
  Data: observational only
\item
  Metric: accuracy, AUC, MSE
\item
  Assumes i.i.d. and stable world
\item
  Loss if wrong: miscalibration
\item
  No feedback with the world
\end{itemize}
\end{column}

\begin{column}{0.5\textwidth}
\textbf{\textcolor{darkgreen}{Decision}}

\begin{itemize}
\tightlist
\item
  Target: \(\arg\max_a \EE[U(Y) | \operatorname{do}(a)]\)
\item
  Data: obs. \(+\) experimental
\item
  Metric: regret, expected utility
\item
  Models the intervention itself
\item
  Loss if wrong: policy harm
\item
  Actions \emph{change} the world
\end{itemize}
\end{column}
\end{columns}

\begin{itemize}
\tightlist
\item
  \textbf{\textcolor{blue}{Key takeaway}}: prediction and
  decision-making live in different mathematical settings; collapsing
  them is a common source of costly errors
\end{itemize}
\end{frame}

\begin{frame}{}
\phantomsection\label{section-2}
\begin{itemize}
\tightlist
\item
  Why Prediction Is Not Enough

  \begin{itemize}
  \tightlist
  \item
    Prediction Pipelines vs.~Decision Pipelines
  \item
    \emph{\textbf{\textcolor{red}{When Prediction Fails}}}
  \item
    Causal Models as the Foundation for Decisions
  \end{itemize}
\item
  Foundations of Decision Theory
\item
  Decision Support with Causal Models
\item
  Sequential Decision-Making and Active Learning
\item
  Uncertainty in Causal Decisions
\end{itemize}
\end{frame}

\begin{frame}{Simpson's Paradox}
\phantomsection\label{simpsons-paradox}
\begin{itemize}
\item
  \textbf{\textcolor{red}{Simpson's paradox}}: a trend in subgroups
  reverses when groups are aggregated

  \begin{itemize}
  \tightlist
  \item
    Cause: a \textbf{\textcolor{darkgreen}{confounder}} is unevenly
    distributed across groups
  \item
    Naive aggregation \(\to\) wrong decision
  \end{itemize}
\item
  \textbf{\textcolor{blue}{Example}}: drug recovery rates
  \begingroup \scriptsize \textbar{} \textbf{Group} \textbar{}
  \textbf{Drug} (recovery) \textbar{} \textbf{No Drug} (recovery)
  \textbar{} \textbar{} --------------- \textbar{} -------------------
  \textbar{} ---------------------- \textbar{} \textbar{} Men \textbar{}
  81 / 87 = 93\% \textbar{} 234 / 270 = 87\% \textbar{} \textbar{} Women
  \textbar{} 192 / 263 = 73\% \textbar{} 55 / 80 = 69\% \textbar{}
  \textbar{} \textbf{Combined} \textbar{} 273 / 350 = 78\% \textbar{}
  289 / 350 = 83\% \textbar{} \endgroup
\item
  Within each subgroup, the drug helps; pooled, it appears to hurt

  \begin{itemize}
  \tightlist
  \item
    Sex is a confounder correlated with both treatment assignment and
    recovery baseline
  \item
    The right decision depends on what we \emph{do} with the drug, not
    what we \emph{see}
  \end{itemize}
\end{itemize}
\end{frame}

\begin{frame}{Simpson's Paradox: Causal Resolution}
\phantomsection\label{simpsons-paradox-causal-resolution}
\begin{itemize}
\tightlist
\item
  The numbers alone cannot tell us whether to prescribe the drug --- we
  need a \textbf{\textcolor{red}{causal graph}}
\end{itemize}

\begin{columns}[T]
\begin{column}{0.55\textwidth}
\begin{itemize}
\tightlist
\item
  \textbf{\textcolor{orange}{Case A}}: Sex \(\to\) Drug, Sex \(\to\)
  Recovery

  \begin{itemize}
  \tightlist
  \item
    Sex confounds treatment choice
  \item
    Correct answer: look at the \textbf{\textcolor{darkgreen}{subgroup}}
    table
  \item
    The drug helps
  \end{itemize}
\item
  \textbf{\textcolor{teal}{Case B}}: Drug \(\to\) Sex (post-treatment),
  Sex \(\to\) Recovery

  \begin{itemize}
  \tightlist
  \item
    Sex is a mediator
  \item
    Correct answer: look at the \textbf{\textcolor{cyan}{aggregate}}
    table
  \item
    The drug hurts
  \end{itemize}
\item
  \emph{Same data, opposite decisions} depending on the causal graph
\end{itemize}
\end{column}

\begin{column}{0.45\textwidth}
\includegraphics{tmp.notes_to_pdf.preprocess_notes.txt.figs/tmp.notes_to_pdf.render_image.1.png}
\end{column}
\end{columns}
\end{frame}

\begin{frame}{Confounding: The Root Cause}
\phantomsection\label{confounding-the-root-cause}
\begin{itemize}
\item
  A \textbf{\textcolor{red}{confounder}} \(Z\) is a common cause of both
  the proposed action \(X\) and the outcome \(Y\)

  \begin{itemize}
  \tightlist
  \item
    Induces a non-causal association \(X \leftrightarrow Y\) through
    \(Z\)
  \item
    Observational \(\Pr(Y | X) \neq \Pr(Y | \operatorname{do}(X))\)
  \end{itemize}
\item
  Without adjusting for \(Z\), any purely predictive model will encode
  the spurious association
\item
  \textbf{\textcolor{darkgreen}{Example}}: ice cream sales and drowning
  deaths are highly correlated

  \begin{itemize}
  \tightlist
  \item
    Predictor: \emph{``more ice cream \(\to\) predict more drownings''}
    (accurate)
  \item
    Decision: \emph{``ban ice cream to reduce drownings''} (disastrous)
  \item
    Confounder: hot weather causes both
  \end{itemize}
\item
  \textbf{\textcolor{blue}{Q}}: Can we avoid this by collecting more
  features?

  \begin{itemize}
  \tightlist
  \item
    Only if the missing confounder is among them \emph{and} properly
    adjusted for
  \end{itemize}
\end{itemize}
\end{frame}

\begin{frame}{Policy Reversal: When Predictions Misguide Action}
\phantomsection\label{policy-reversal-when-predictions-misguide-action}
\begin{itemize}
\item
  A \textbf{\textcolor{red}{policy reversal}} is when an optimal-looking
  decision based on predictions becomes harmful once enacted

  \begin{itemize}
  \tightlist
  \item
    Common in medicine, economics, operations research
  \end{itemize}
\item
  \textbf{\textcolor{orange}{Mechanism}}: acting on \(X\) breaks the
  correlation with the unobserved confounder \(Z\), so the observed
  association no longer holds
\item
  \textbf{\textcolor{darkgreen}{Example}}: emergency room triage

  \begin{itemize}
  \tightlist
  \item
    Predictor: patients with asthma + pneumonia have \emph{lower}
    mortality
  \item
    Naive policy: send them home first (they're low-risk)
  \item
    Reality: they had lower mortality \emph{because} they received
    aggressive care

    \begin{itemize}
    \tightlist
    \item
      Removing the care removes the protection
    \end{itemize}
  \end{itemize}
\item
  \textbf{\textcolor{teal}{Lesson}}: predictive patterns in historical
  data reflect the \textbf{policy that generated the data}, not what
  happens when we change it
\end{itemize}
\end{frame}

\begin{frame}{}
\phantomsection\label{section-3}
\begin{itemize}
\tightlist
\item
  Why Prediction Is Not Enough

  \begin{itemize}
  \tightlist
  \item
    Prediction Pipelines vs.~Decision Pipelines
  \item
    When Prediction Fails
  \item
    \emph{\textbf{\textcolor{red}{Causal Models as the Foundation for Decisions}}}
  \end{itemize}
\item
  Foundations of Decision Theory
\item
  Decision Support with Causal Models
\item
  Sequential Decision-Making and Active Learning
\item
  Uncertainty in Causal Decisions
\end{itemize}
\end{frame}

\begin{frame}{Causal Models for Decisions}
\phantomsection\label{causal-models-for-decisions}
\begin{itemize}
\item
  A \textbf{\textcolor{red}{causal model}} specifies, for each variable
  \(X_i\):

  \begin{itemize}
  \tightlist
  \item
    Its direct causes (parents) \(\text{Pa}(X_i)\)
  \item
    A structural mechanism \(X_i = f_i(\text{Pa}(X_i), U_i)\) with noise
    \(U_i\)
  \end{itemize}
\item
  Enables reasoning about:

  \begin{itemize}
  \tightlist
  \item
    \textbf{Observation}: \(\Pr(Y | X = x)\) (seeing \(x\))
  \item
    \textbf{Intervention}: \(\Pr(Y | \operatorname{do}(X = x))\)
    (setting \(x\))
  \item
    \textbf{Counterfactual}: \(\Pr(Y_{X=x} | X = x', Y = y')\) (had
    \(x\) been different)
  \end{itemize}
\item
  For decision-making we primarily need the
  \textbf{\textcolor{darkgreen}{interventional}} layer
\item
  \textbf{\textcolor{blue}{Example}}: a marketing team can observe that
  users who open an email buy more; only a causal model can say whether
  \emph{sending} the email causes purchases
\end{itemize}
\end{frame}

\begin{frame}{Decision Loop with Causal Models}
\phantomsection\label{decision-loop-with-causal-models}
\includegraphics{tmp.notes_to_pdf.preprocess_notes.txt.figs/tmp.notes_to_pdf.render_image.2.png}

\begin{itemize}
\tightlist
\item
  \textbf{\textcolor{red}{Key insight}}: the causal model sits between
  data and decisions

  \begin{itemize}
  \tightlist
  \item
    Pure prediction skips the middle step
  \item
    New data from interventions updates the model, closing the loop
  \end{itemize}
\end{itemize}
\end{frame}

\begin{frame}{}
\phantomsection\label{section-4}
\begin{itemize}
\tightlist
\item
  Why Prediction Is Not Enough
\item
  \emph{\textbf{\textcolor{red}{Foundations of Decision Theory}}}

  \begin{itemize}
  \tightlist
  \item
    Utility Functions and Expected Utility
  \item
    Causal Interventions and Expected Outcomes
  \item
    Risk Preferences and Multi-Criteria Trade-offs
  \end{itemize}
\item
  Decision Support with Causal Models
\item
  Sequential Decision-Making and Active Learning
\item
  Uncertainty in Causal Decisions
\end{itemize}
\end{frame}

\begin{frame}{}
\phantomsection\label{section-5}
\begin{itemize}
\tightlist
\item
  Why Prediction Is Not Enough
\item
  Foundations of Decision Theory

  \begin{itemize}
  \tightlist
  \item
    \emph{\textbf{\textcolor{red}{Utility Functions and Expected Utility}}}
  \item
    Causal Interventions and Expected Outcomes
  \item
    Risk Preferences and Multi-Criteria Trade-offs
  \end{itemize}
\item
  Decision Support with Causal Models
\item
  Sequential Decision-Making and Active Learning
\item
  Uncertainty in Causal Decisions
\end{itemize}
\end{frame}

\begin{frame}{Utility Functions}
\phantomsection\label{utility-functions}
\begin{itemize}
\tightlist
\item
  A \textbf{\textcolor{red}{utility function}}
  \(U: \mathcal{Y} \to \RR\) maps each outcome \(y\) to a real-valued
  preference score

  \begin{itemize}
  \tightlist
  \item
    Higher is better
  \item
    Encodes the decision-maker's goals on a single scale
  \end{itemize}
\item
  \textbf{\textcolor{darkgreen}{Properties}}:

  \begin{itemize}
  \tightlist
  \item
    Utilities are \emph{ordinal}: rank is what matters, not absolute
    values
  \item
    Utilities are \emph{subjective}: different stakeholders can have
    different \(U\)
  \item
    Defined up to positive affine transformations \(U \to a U + b\),
    \(a > 0\)
  \end{itemize}
\item
  \textbf{\textcolor{blue}{Example}}: a hospital utility

  \begin{itemize}
  \tightlist
  \item
    Full recovery: \(U = 1.0\)
  \item
    Recovery with complications: \(U = 0.7\)
  \item
    Death: \(U = 0.0\)
  \item
    A \(10\%\) complication reduction is worth \(0.03\) in utility on
    average
  \end{itemize}
\end{itemize}
\end{frame}

\begin{frame}{Expected Utility Principle}
\phantomsection\label{expected-utility-principle}
\begin{itemize}
\item
  Under uncertainty, a rational agent chooses the action \(a\) that
  maximizes \textbf{\textcolor{red}{expected utility}}: \[
  a^* = \arg\max_{a \in \mathcal{A}} \EE_{Y \sim \Pr(\cdot | a)}[U(Y)]
  \]
\item
  \textbf{\textcolor{darkgreen}{Von Neumann-Morgenstern theorem (1944)}}:
  if preferences over lotteries satisfy 4 axioms (completeness,
  transitivity, continuity, independence), there exists a utility
  function \(U\) such that the agent acts \emph{as if} maximizing
  expected utility
\item
  \textbf{\textcolor{blue}{Example}}: buying insurance

  \begin{itemize}
  \tightlist
  \item
    Premium \(c\) for certain
  \item
    Uninsured loss \(L\) with probability \(p\)
  \item
    Insure if \(U(W - c) > p \cdot U(W - L) + (1 - p) \cdot U(W)\)
  \item
    Risk-averse utilities make this inequality easier to satisfy
  \end{itemize}
\end{itemize}
\end{frame}

\begin{frame}{Expected Utility: A Simple Example}
\phantomsection\label{expected-utility-a-simple-example}
\begin{itemize}
\item
  \textbf{\textcolor{red}{Setup}}: a doctor chooses between two
  treatments

  \begin{itemize}
  \tightlist
  \item
    Treatment 1: cures with probability \(0.9\), no side effect
  \item
    Treatment 2: cures with probability \(0.95\), mild side effect
    (\(U = 0.8\))
  \end{itemize}
\item
  \textbf{\textcolor{darkgreen}{Utility}}:

  \begin{itemize}
  \tightlist
  \item
    Cure with no side effect: \(1.0\)
  \item
    Cure with side effect: \(0.8\)
  \item
    No cure: \(0.0\)
  \end{itemize}
\item
  \textbf{\textcolor{blue}{Computation}}: \begin{align*}
  \EE[U | T_1] &= 0.9 \cdot 1.0 + 0.1 \cdot 0.0 = 0.90 \\
  \EE[U | T_2] &= 0.95 \cdot 0.8 + 0.05 \cdot 0.0 = 0.76
  \end{align*}
\item
  \textbf{\textcolor{violet}{Decision}}: choose \(T_1\)

  \begin{itemize}
  \tightlist
  \item
    Even though \(T_2\) is more effective, its side-effect penalty
    outweighs the gain
  \end{itemize}
\end{itemize}
\end{frame}

\begin{frame}{}
\phantomsection\label{section-6}
\begin{itemize}
\tightlist
\item
  Why Prediction Is Not Enough
\item
  Foundations of Decision Theory

  \begin{itemize}
  \tightlist
  \item
    Utility Functions and Expected Utility
  \item
    \emph{\textbf{\textcolor{red}{Causal Interventions and Expected Outcomes}}}
  \item
    Risk Preferences and Multi-Criteria Trade-offs
  \end{itemize}
\item
  Decision Support with Causal Models
\item
  Sequential Decision-Making and Active Learning
\item
  Uncertainty in Causal Decisions
\end{itemize}
\end{frame}

\begin{frame}{From Observational to Interventional Expectation}
\phantomsection\label{from-observational-to-interventional-expectation}
\begin{itemize}
\item
  Classical decision theory uses \(\Pr(Y | a)\) --- treated as the
  distribution of \(Y\) given action \(a\)

  \begin{itemize}
  \tightlist
  \item
    If the data is observational, \(\Pr(Y | a)\) includes confounding
    bias
  \item
    Plugging it into expected utility yields a
    \textbf{\textcolor{red}{biased}} action ranking
  \end{itemize}
\item
  Causal decision theory replaces it with the
  \textbf{\textcolor{darkgreen}{interventional distribution}}: \[
  a^* = \arg\max_{a \in \mathcal{A}} \EE_{Y \sim \Pr(Y | \operatorname{do}(a))}[U(Y)]
  \]
\item
  \textbf{\textcolor{blue}{Equivalent under}}:

  \begin{itemize}
  \tightlist
  \item
    Randomized experiments (no confounders)
  \item
    Successful back-door or front-door adjustment
  \item
    Complete knowledge of the causal graph + sufficient data
  \end{itemize}
\end{itemize}
\end{frame}

\begin{frame}{Intervention in an Influence Graph}
\phantomsection\label{intervention-in-an-influence-graph}
\begin{columns}[T]
\begin{column}{0.55\textwidth}
\begin{itemize}
\item
  \textbf{\textcolor{red}{Do-operator}} \(\operatorname{do}(X = x)\)
  mutates the causal graph:

  \begin{itemize}
  \tightlist
  \item
    Remove all edges \emph{into} \(X\)
  \item
    Fix \(X\) to the value \(x\)
  \item
    Propagate through the remaining structural equations
  \end{itemize}
\item
  Turns a \emph{passive} question into an \emph{active} one
\item
  \textbf{\textcolor{blue}{Example}}:
  \(\operatorname{do}(\text{Drug} = 1)\)

  \begin{itemize}
  \tightlist
  \item
    Drug is no longer chosen by the physician
  \item
    All causal arrows from Sex, Age, etc. \emph{into} Drug are cut
  \item
    Resulting distribution reflects a policy, not the observed world
  \end{itemize}
\end{itemize}
\end{column}

\begin{column}{0.45\textwidth}
\includegraphics{tmp.notes_to_pdf.preprocess_notes.txt.figs/tmp.notes_to_pdf.render_image.3.png}

\scriptsize

\emph{Edge \(Z \to X\) is cut under} \(\operatorname{do}(X)\)
\end{column}
\end{columns}
\end{frame}

\begin{frame}{Expected Outcome of an Intervention}
\phantomsection\label{expected-outcome-of-an-intervention}
\begin{itemize}
\item
  Given a causal graph and estimated mechanisms, compute the expected
  outcome under each candidate action: \[
  \EE[Y | \operatorname{do}(a)] = \sum_{y} y \cdot \Pr(Y = y | \operatorname{do}(A = a))
  \]
\item
  \textbf{\textcolor{red}{Adjustment formula}} (back-door): \[
  \Pr(Y | \operatorname{do}(A)) = \sum_{Z} \Pr(Y | A, Z) \Pr(Z)
  \] where \(Z\) is a back-door adjustment set
\item
  \textbf{\textcolor{blue}{Example}}: drug decision

  \begin{itemize}
  \tightlist
  \item
    Observational: \(\Pr(\text{recovery} | \text{drug}) = 0.78\)
  \item
    Adjusted: \(\sum_z \Pr(\text{rec} | \text{drug}, z) \Pr(z) = 0.83\)
  \item
    The adjusted value is what plugs into expected utility
  \end{itemize}
\end{itemize}
\end{frame}

\begin{frame}{}
\phantomsection\label{section-7}
\begin{itemize}
\tightlist
\item
  Why Prediction Is Not Enough
\item
  Foundations of Decision Theory

  \begin{itemize}
  \tightlist
  \item
    Utility Functions and Expected Utility
  \item
    Causal Interventions and Expected Outcomes
  \item
    \emph{\textbf{\textcolor{red}{Risk Preferences and Multi-Criteria Trade-offs}}}
  \end{itemize}
\item
  Decision Support with Causal Models
\item
  Sequential Decision-Making and Active Learning
\item
  Uncertainty in Causal Decisions
\end{itemize}
\end{frame}

\begin{frame}{Risk Preferences}
\phantomsection\label{risk-preferences}
\begin{itemize}
\tightlist
\item
  Expected utility handles risk through the \emph{shape} of \(U\):

  \begin{itemize}
  \tightlist
  \item
    \textbf{\textcolor{red}{Risk-neutral}}: \(U\) is linear,
    \(U(y) = y\) --- care only about the mean
  \item
    \textbf{\textcolor{orange}{Risk-averse}}: \(U\) is concave ---
    penalize variance
  \item
    \textbf{\textcolor{darkgreen}{Risk-seeking}}: \(U\) is convex ---
    reward variance
  \end{itemize}
\item
  \textbf{\textcolor{teal}{Certainty equivalent}} \(y_{CE}\): the sure
  outcome the agent values as much as a risky prospect

  \begin{itemize}
  \tightlist
  \item
    Risk-averse: \(y_{CE} < \EE[Y]\)
  \item
    Gap \(\EE[Y] - y_{CE}\) is the
    \textbf{\textcolor{cyan}{risk premium}}
  \end{itemize}
\item
  \textbf{\textcolor{blue}{Example}}: \(U(y) = \log(y)\)

  \begin{itemize}
  \tightlist
  \item
    A 50/50 bet between \$100 and \$400 has mean \$250
  \item
    \(y_{CE} = \exp(0.5 \log 100 + 0.5 \log 400) = \$200 < \$250\)
  \end{itemize}
\end{itemize}
\end{frame}

\begin{frame}{Visualizing Risk Aversion}
\phantomsection\label{visualizing-risk-aversion}
\includegraphics{tmp.notes_to_pdf.preprocess_notes.txt.figs/tmp.notes_to_pdf.render_image.4.png}

\begin{itemize}
\tightlist
\item
  \textbf{\textcolor{red}{Key takeaway}}: a concave \(U\) trades
  expected return for reduced variance

  \begin{itemize}
  \tightlist
  \item
    Choose actions whose outcome distribution is \emph{both} high in
    mean and narrow in spread
  \end{itemize}
\end{itemize}
\end{frame}

\begin{frame}{Multi-Criteria Trade-offs}
\phantomsection\label{multi-criteria-trade-offs}
\begin{itemize}
\tightlist
\item
  Real decisions weigh \textbf{\textcolor{red}{multiple objectives}}
  that pull in different directions

  \begin{itemize}
  \tightlist
  \item
    Clinical: efficacy, side effects, cost, patient burden
  \item
    Business: revenue, churn, brand, regulatory risk
  \end{itemize}
\item
  \textbf{\textcolor{darkgreen}{Scalarization}}: collapse into a single
  utility \[
  U(y) = \sum_{k} w_k \cdot u_k(y_k)
  \]

  \begin{itemize}
  \tightlist
  \item
    Weights \(w_k\) elicited from stakeholders
  \item
    Works when trade-offs are approximately linear
  \end{itemize}
\item
  \textbf{\textcolor{blue}{Pareto frontier}}: keep only non-dominated
  actions

  \begin{itemize}
  \tightlist
  \item
    Avoid picking a single \(U\) until preferences are clear
  \item
    Present trade-offs to the decision-maker visually
  \end{itemize}
\item
  \textbf{\textcolor{violet}{Example}}: recommending a cancer therapy

  \begin{itemize}
  \tightlist
  \item
    Objectives: 5-year survival, quality of life, cost
  \item
    Elicit weights; or present the Pareto front of treatments
  \end{itemize}
\end{itemize}
\end{frame}

\begin{frame}{Risk and Trade-offs: Practical Advice}
\phantomsection\label{risk-and-trade-offs-practical-advice}
\begin{itemize}
\tightlist
\item
  \textbf{\textcolor{red}{Pros}}

  \begin{itemize}
  \tightlist
  \item
    Expected utility unifies risk, uncertainty, and multi-criteria
    reasoning
  \item
    Calibrated utilities make implicit trade-offs explicit
  \item
    Maps cleanly onto Bayesian posteriors
  \end{itemize}
\item
  \textbf{\textcolor{darkgreen}{Cons}}

  \begin{itemize}
  \tightlist
  \item
    Utility elicitation is hard --- stakeholders struggle to quantify
    preferences
  \item
    Sensitivity: small changes in \(U\) can flip the optimal action
  \item
    Behavioral humans \emph{systematically} violate the axioms (Kahneman
    \& Tversky)
  \end{itemize}
\item
  \textbf{\textcolor{blue}{Rule of thumb}}: perform a
  \textbf{\textcolor{violet}{sensitivity analysis}} over \(U\) --- if
  the optimal action is stable across a plausible range of utilities,
  the decision is robust
\end{itemize}
\end{frame}

\begin{frame}{}
\phantomsection\label{section-8}
\begin{itemize}
\tightlist
\item
  Why Prediction Is Not Enough
\item
  Foundations of Decision Theory
\item
  \emph{\textbf{\textcolor{red}{Decision Support with Causal Models}}}

  \begin{itemize}
  \tightlist
  \item
    Influence Diagrams
  \item
    Bayesian Decision-Making
  \item
    Prior Elicitation
  \end{itemize}
\item
  Sequential Decision-Making and Active Learning
\item
  Uncertainty in Causal Decisions
\end{itemize}
\end{frame}

\begin{frame}{}
\phantomsection\label{section-9}
\begin{itemize}
\tightlist
\item
  Why Prediction Is Not Enough
\item
  Foundations of Decision Theory
\item
  Decision Support with Causal Models

  \begin{itemize}
  \tightlist
  \item
    \emph{\textbf{\textcolor{red}{Influence Diagrams}}}
  \item
    Bayesian Decision-Making
  \item
    Prior Elicitation
  \end{itemize}
\item
  Sequential Decision-Making and Active Learning
\item
  Uncertainty in Causal Decisions
\end{itemize}
\end{frame}

\begin{frame}{Influence Diagrams: Motivation}
\phantomsection\label{influence-diagrams-motivation}
\begin{itemize}
\item
  Bayesian networks model \emph{uncertainty}; we need to also represent
  \emph{decisions} and \emph{preferences}
\item
  An \textbf{\textcolor{red}{influence diagram}} extends a causal DAG
  with three kinds of nodes:

  \begin{itemize}
  \tightlist
  \item
    \textbf{\textcolor{orange}{Chance nodes}} (ovals): random variables,
    same as in a Bayesian network
  \item
    \textbf{\textcolor{darkgreen}{Decision nodes}} (rectangles):
    variables set by the decision-maker
  \item
    \textbf{\textcolor{teal}{Utility nodes}} (diamonds): deterministic
    functions yielding utility
  \end{itemize}
\item
  \textbf{\textcolor{cyan}{Purpose}}: a compact, graphical specification
  of a decision problem that can be solved by message passing
\item
  \textbf{\textcolor{blue}{Example}}: whether to run a diagnostic test
  before treating a patient

  \begin{itemize}
  \tightlist
  \item
    Chance: disease status, test result, outcome
  \item
    Decision: order test?, choose therapy?
  \item
    Utility: cost of test, value of outcome
  \end{itemize}
\end{itemize}
\end{frame}

\begin{frame}{Influence Diagram: Anatomy}
\phantomsection\label{influence-diagram-anatomy}
\includegraphics{tmp.notes_to_pdf.preprocess_notes.txt.figs/tmp.notes_to_pdf.render_image.5.png}

\begin{itemize}
\tightlist
\item
  Arrows into decision nodes denote
  \textbf{\textcolor{red}{information available}} at decision time (no
  causal meaning)
\item
  Arrows into chance and utility nodes retain their
  \textbf{\textcolor{blue}{causal}} interpretation
\end{itemize}
\end{frame}

\begin{frame}{Solving an Influence Diagram}
\phantomsection\label{solving-an-influence-diagram}
\begin{itemize}
\item
  \textbf{\textcolor{red}{Input}}: influence diagram with CPTs for
  chance nodes and a utility table
\item
  \textbf{\textcolor{orange}{Output}}: optimal policy \(\pi^*\) =
  function from information to actions
\item
  \textbf{\textcolor{darkgreen}{Steps}} (variable elimination):

  \begin{enumerate}
  \tightlist
  \item
    Eliminate chance nodes by summing them out
  \item
    Eliminate decision nodes by \textbf{\textcolor{teal}{maximizing}}
    over their domain
  \item
    Process in reverse topological order
  \item
    Record the arg-max at each decision node as the optimal policy
  \end{enumerate}
\item
  \textbf{\textcolor{cyan}{Complexity}}:

  \begin{itemize}
  \tightlist
  \item
    Time: \(O(|\mathcal{A}|^d \cdot \prod_i |X_i|)\) in the worst case
  \item
    Tractable when the graph has low treewidth
  \end{itemize}
\item
  \textbf{\textcolor{blue}{Key insight}}: the same graphical machinery
  used for Bayesian inference also supports optimal decision-making
\end{itemize}
\end{frame}

\begin{frame}{Influence Diagrams: Pros and Cons}
\phantomsection\label{influence-diagrams-pros-and-cons}
\begin{itemize}
\tightlist
\item
  \textbf{\textcolor{red}{Pros}}

  \begin{itemize}
  \tightlist
  \item
    Explicit separation of beliefs, choices, and preferences
  \item
    Easy to communicate to non-technical stakeholders
  \item
    Supports value of information analysis out of the box
  \end{itemize}
\item
  \textbf{\textcolor{darkgreen}{Cons}}

  \begin{itemize}
  \tightlist
  \item
    Requires fully specified CPTs and utilities --- often heroic
  \item
    Exact solution is intractable for large continuous problems
  \item
    Assumes a single rational decision-maker with stable utilities
  \end{itemize}
\item
  \textbf{\textcolor{blue}{Example}}: the oil wildcatter problem
  (classic)

  \begin{itemize}
  \tightlist
  \item
    Decide whether to drill, after optionally running a seismic test
  \item
    Closed-form optimal policy using a 4-node influence diagram
  \end{itemize}
\end{itemize}
\end{frame}

\begin{frame}{}
\phantomsection\label{section-10}
\begin{itemize}
\tightlist
\item
  Why Prediction Is Not Enough
\item
  Foundations of Decision Theory
\item
  Decision Support with Causal Models

  \begin{itemize}
  \tightlist
  \item
    Influence Diagrams
  \item
    \emph{\textbf{\textcolor{red}{Bayesian Decision-Making}}}
  \item
    Prior Elicitation
  \end{itemize}
\item
  Sequential Decision-Making and Active Learning
\item
  Uncertainty in Causal Decisions
\end{itemize}
\end{frame}

\begin{frame}{Posteriors to Optimal Actions}
\phantomsection\label{posteriors-to-optimal-actions}
\begin{itemize}
\item
  Bayesian decision-making composes two steps:

  \begin{enumerate}
  \tightlist
  \item
    \textbf{\textcolor{red}{Inference}}: compute the posterior
    \(\Pr(\theta | \mathcal{D})\) from data
  \item
    \textbf{\textcolor{orange}{Decision}}: choose the action that
    maximizes \emph{posterior} expected utility \[
    a^* = \arg\max_{a \in \mathcal{A}} \EE_{\theta \sim \Pr(\theta | \mathcal{D})}
       \left[ \EE_{Y \sim \Pr(Y | \operatorname{do}(a), \theta)}[U(Y)] \right]
    \]
  \end{enumerate}
\item
  The \textbf{\textcolor{darkgreen}{inner expectation}} is over outcome
  uncertainty given \(\theta\) (aleatoric)
\item
  The \textbf{\textcolor{teal}{outer expectation}} is over parameter
  uncertainty (epistemic)
\item
  \textbf{\textcolor{cyan}{Key property}}: Bayesian decision rule is
  \emph{admissible} --- no other rule dominates it under all priors
\end{itemize}
\end{frame}

\begin{frame}{Bayesian Decision Rule: Example}
\phantomsection\label{bayesian-decision-rule-example}
\begin{itemize}
\item
  \textbf{\textcolor{red}{Scenario}}: A/B test of a new checkout flow

  \begin{itemize}
  \tightlist
  \item
    Parameter \(\theta\): conversion rate lift of variant B over A
  \item
    Prior: \(\theta \sim \mathcal{N}(0, 0.02^2)\)
  \item
    Data \(\mathcal{D}\): \(n = 5000\) users per arm; observed lift
    \(0.01\)
  \end{itemize}
\item
  \textbf{Posterior}:
  \(\theta | \mathcal{D} \sim \mathcal{N}(0.008, 0.004^2)\)

  \begin{itemize}
  \tightlist
  \item
    \(\Pr(\theta > 0 | \mathcal{D}) \approx 0.977\)
  \end{itemize}
\item
  \textbf{\textcolor{darkgreen}{Utility of launching}}
  \(U(\theta) = \theta - c\), where \(c = 0.002\) represents the cost of
  maintaining the new code
\item
  \textbf{\textcolor{blue}{Expected utility}}:

  \begin{itemize}
  \tightlist
  \item
    \(\EE[U | \text{launch}] = 0.008 - 0.002 = 0.006 > 0\)
  \item
    \(\EE[U | \text{keep A}] = 0\)
  \item
    \textbf{\textcolor{violet}{Decision}}: launch B
  \end{itemize}
\end{itemize}
\end{frame}

\begin{frame}{Bayesian vs.~Frequentist Decisions}
\phantomsection\label{bayesian-vs.-frequentist-decisions}
\begin{columns}[T]
\begin{column}{0.5\textwidth}
\textbf{\textcolor{red}{Bayesian}}

\begin{itemize}
\item
  Posterior \(\Pr(\theta | \mathcal{D})\)
\item
  Acts on full distribution
\item
  Incorporates prior beliefs
\item
  Natural for sequential updating
\item
  Connects directly to utility
\item
  \textbf{\textcolor{orange}{Pros}}

  \begin{itemize}
  \tightlist
  \item
    Coherent under uncertainty
  \item
    Works for small samples
  \end{itemize}
\item
  \textbf{\textcolor{darkgreen}{Cons}}

  \begin{itemize}
  \tightlist
  \item
    Requires priors
  \item
    Computationally heavier
  \end{itemize}
\end{itemize}
\end{column}

\begin{column}{0.5\textwidth}
\textbf{\textcolor{teal}{Frequentist}}

\begin{itemize}
\item
  Point estimate \(\hat{\theta}\) + CI
\item
  Acts on test outcomes
\item
  No priors
\item
  Decision via hypothesis test
\item
  Utility often implicit
\item
  \textbf{\textcolor{cyan}{Pros}}

  \begin{itemize}
  \tightlist
  \item
    Fewer subjective inputs
  \item
    Standardized in many fields
  \end{itemize}
\item
  \textbf{\textcolor{blue}{Cons}}

  \begin{itemize}
  \tightlist
  \item
    Thresholds (e.g., \(p < 0.05\)) are not utility-optimal
  \item
    Harder to combine with new data
  \end{itemize}
\end{itemize}
\end{column}
\end{columns}

\begin{itemize}
\tightlist
\item
  \textbf{\textcolor{violet}{Bottom line}}: Bayesian framing slots
  naturally into causal decision-making because both use joint
  distributions end-to-end
\end{itemize}
\end{frame}

\begin{frame}{}
\phantomsection\label{section-11}
\begin{itemize}
\tightlist
\item
  Why Prediction Is Not Enough
\item
  Foundations of Decision Theory
\item
  Decision Support with Causal Models

  \begin{itemize}
  \tightlist
  \item
    Influence Diagrams
  \item
    Bayesian Decision-Making
  \item
    \emph{\textbf{\textcolor{red}{Prior Elicitation}}}
  \end{itemize}
\item
  Sequential Decision-Making and Active Learning
\item
  Uncertainty in Causal Decisions
\end{itemize}
\end{frame}

\begin{frame}{Prior Elicitation: What and Why}
\phantomsection\label{prior-elicitation-what-and-why}
\begin{itemize}
\tightlist
\item
  In causal decision problems, the prior encodes beliefs about the
  \textbf{causal effect} \(\theta\) before seeing data

  \begin{itemize}
  \tightlist
  \item
    Size, sign, and uncertainty of \(\theta\)
  \end{itemize}
\item
  Prior elicitation is especially important when:

  \begin{itemize}
  \tightlist
  \item
    Data is scarce or expensive (medicine, policy)
  \item
    Experiments are unethical or impossible
  \item
    Stakeholders disagree
  \end{itemize}
\item
  \textbf{\textcolor{red}{Q}}: Where do priors come from?

  \begin{itemize}
  \tightlist
  \item
    Literature meta-analyses
  \item
    Expert judgment
  \item
    Past experiments or pilot studies
  \item
    Regulatory or legal constraints
  \end{itemize}
\end{itemize}
\end{frame}

\begin{frame}{Methods for Eliciting Priors}
\phantomsection\label{methods-for-eliciting-priors}
\begin{itemize}
\item
  \textbf{\textcolor{red}{Quantile elicitation}}: ask experts for
  10/50/90 percentiles of \(\theta\), then fit a distribution

  \begin{itemize}
  \tightlist
  \item
    E.g., ``10\% chance lift is below \(0.5\%\), 50\% below \(1\%\),
    90\% below \(3\%\)''
  \end{itemize}
\item
  \textbf{\textcolor{orange}{Pairwise judgments}}: elicit relative
  probabilities of outcomes, combine via a log-linear model
\item
  \textbf{\textcolor{darkgreen}{Historical calibration}}: compare past
  predictions against actuals to adjust for bias and overconfidence
\item
  \textbf{\textcolor{teal}{Reference priors}}: use weakly informative
  priors (Jeffreys, unit information) when knowledge is vague

  \begin{itemize}
  \tightlist
  \item
    Avoid flat priors in high dimensions --- they are not uninformative
  \end{itemize}
\item
  \textbf{\textcolor{cyan}{Example}}: a clinical trial team estimates a
  20\% prior probability that a new drug is at least as effective as
  standard of care, based on meta-analysis of 14 similar compounds
\end{itemize}
\end{frame}

\begin{frame}{Pitfalls and Best Practices}
\phantomsection\label{pitfalls-and-best-practices}
\begin{itemize}
\tightlist
\item
  \textbf{\textcolor{red}{Pitfalls}}

  \begin{itemize}
  \tightlist
  \item
    \textbf{\textcolor{orange}{Anchoring}}: experts cling to a
    first-mentioned number
  \item
    \textbf{\textcolor{darkgreen}{Overconfidence}}: intervals are
    consistently too narrow
  \item
    \textbf{\textcolor{teal}{Framing effects}}: elicit the \emph{same}
    quantity two ways --- answers differ
  \item
    \textbf{\textcolor{cyan}{Confusing}} aleatoric vs.~epistemic
    uncertainty --- see later section
  \end{itemize}
\item
  \textbf{\textcolor{blue}{Best practices}}

  \begin{itemize}
  \tightlist
  \item
    Elicit from multiple experts independently, then aggregate
  \item
    Cross-validate with subsets of historical data
  \item
    Perform \textbf{\textcolor{violet}{prior sensitivity analysis}}:
    does the decision change under plausible alternative priors?
  \item
    Document the elicitation process --- decisions may need defending
  \end{itemize}
\end{itemize}
\end{frame}

\begin{frame}{}
\phantomsection\label{section-12}
\begin{itemize}
\tightlist
\item
  Why Prediction Is Not Enough
\item
  Foundations of Decision Theory
\item
  Decision Support with Causal Models
\item
  \emph{\textbf{\textcolor{red}{Sequential Decision-Making and Active Learning}}}

  \begin{itemize}
  \tightlist
  \item
    Value of Information
  \item
    Exploration vs.~Exploitation
  \item
    Bayesian Optimization for Experimentation
  \end{itemize}
\item
  Uncertainty in Causal Decisions
\end{itemize}
\end{frame}

\begin{frame}{}
\phantomsection\label{section-13}
\begin{itemize}
\tightlist
\item
  Why Prediction Is Not Enough
\item
  Foundations of Decision Theory
\item
  Decision Support with Causal Models
\item
  Sequential Decision-Making and Active Learning

  \begin{itemize}
  \tightlist
  \item
    \emph{\textbf{\textcolor{red}{Value of Information}}}
  \item
    Exploration vs.~Exploitation
  \item
    Bayesian Optimization for Experimentation
  \end{itemize}
\item
  Uncertainty in Causal Decisions
\end{itemize}
\end{frame}

\begin{frame}{Value of Information: Motivation}
\phantomsection\label{value-of-information-motivation}
\begin{itemize}
\item
  Before committing to a costly action, a decision-maker may gather more
  data

  \begin{itemize}
  \tightlist
  \item
    Diagnostic test, survey, pilot experiment, additional A/B arm
  \end{itemize}
\item
  \textbf{\textcolor{red}{Q}}: Is the extra information worth its cost?
\item
  \textbf{\textcolor{darkgreen}{Value of information (VOI)}} quantifies
  the \emph{expected} gain in utility from obtaining new evidence
  \emph{before} making a decision

  \begin{itemize}
  \tightlist
  \item
    A central concept in decision analysis
  \end{itemize}
\item
  \textbf{\textcolor{blue}{Example}}: a company considers an expensive
  market survey before a \$10M product launch

  \begin{itemize}
  \tightlist
  \item
    Launch without: expected utility \(\EE[U]\)
  \item
    Launch with survey outcome \(E\):
    \(\EE_E[\max_a \EE[U(Y) | \operatorname{do}(a), E]]\)
  \item
    Run the survey only if the difference exceeds its cost
  \end{itemize}
\end{itemize}
\end{frame}

\begin{frame}{Expected Value of Perfect Information (EVPI)}
\phantomsection\label{expected-value-of-perfect-information-evpi}
\begin{itemize}
\item
  \textbf{\textcolor{red}{EVPI}} is the upper bound on what any
  information-gathering activity could be worth: \[
  \text{EVPI} = \EE_{\theta} \left[ \max_a U(a, \theta) \right] -
                \max_a \EE_\theta[U(a, \theta)]
  \]
\item
  Interpret the two terms:

  \begin{itemize}
  \tightlist
  \item
    First: utility if we could learn \(\theta\) \emph{perfectly} before
    deciding
  \item
    Second: utility of the best action under current uncertainty
  \end{itemize}
\item
  \textbf{\textcolor{blue}{Key properties}}:

  \begin{itemize}
  \tightlist
  \item
    \(\text{EVPI} \geq 0\) always
  \item
    \(\text{EVPI} = 0\) iff the current best action is optimal for
    \emph{every} \(\theta\)
  \item
    Any test costing more than EVPI is not worth running, under any
    accuracy
  \end{itemize}
\end{itemize}
\end{frame}

\begin{frame}{Expected Value of Sample Information (EVSI)}
\phantomsection\label{expected-value-of-sample-information-evsi}
\begin{itemize}
\tightlist
\item
  Real information sources are \textbf{\textcolor{red}{imperfect}}

  \begin{itemize}
  \tightlist
  \item
    EVSI relaxes EVPI to consider noisy evidence \(E\)
  \item
    \(\text{EVSI}(E) = \EE_E[\max_a \EE[U | a, E]] - \max_a \EE[U | a]\)
  \end{itemize}
\item
  \(0 \leq \text{EVSI}(E) \leq \text{EVPI}\)

  \begin{itemize}
  \tightlist
  \item
    Noisier evidence \(\to\) smaller EVSI
  \item
    Tighter prior \(\to\) smaller EVSI (less to learn)
  \end{itemize}
\item
  \textbf{\textcolor{darkgreen}{Example}}: diagnostic test with
  sensitivity \(0.9\) and specificity \(0.85\)

  \begin{itemize}
  \tightlist
  \item
    Compute EVSI by integrating over all possible \((E = +, E = -)\)
    outcomes
  \item
    Run the test if \(\text{EVSI} - \text{cost} > 0\)
  \end{itemize}
\item
  \textbf{\textcolor{blue}{Practical tip}}: run a quick EVPI screen
  first --- if EVPI is tiny, skip the EVSI computation entirely
\end{itemize}
\end{frame}

\begin{frame}{}
\phantomsection\label{section-14}
\begin{itemize}
\tightlist
\item
  Why Prediction Is Not Enough
\item
  Foundations of Decision Theory
\item
  Decision Support with Causal Models
\item
  Sequential Decision-Making and Active Learning

  \begin{itemize}
  \tightlist
  \item
    Value of Information
  \item
    \emph{\textbf{\textcolor{red}{Exploration vs. Exploitation}}}
  \item
    Bayesian Optimization for Experimentation
  \end{itemize}
\item
  Uncertainty in Causal Decisions
\end{itemize}
\end{frame}

\begin{frame}{Exploration vs.~Exploitation}
\phantomsection\label{exploration-vs.-exploitation}
\begin{itemize}
\item
  A sequential decision-maker faces a tension:

  \begin{itemize}
  \tightlist
  \item
    \textbf{\textcolor{red}{Exploit}}: pick the action that
    \emph{currently looks best}
  \item
    \textbf{\textcolor{darkgreen}{Explore}}: try other actions to
    \emph{reduce uncertainty} about their effects
  \end{itemize}
\item
  Pure exploitation can lock in a suboptimal choice based on early noise

  \begin{itemize}
  \tightlist
  \item
    Pure exploration wastes samples on clearly bad options
  \end{itemize}
\item
  \textbf{\textcolor{blue}{Example}}: a doctor who always prescribes the
  first drug that worked early in her career will never learn about
  newer, better options
\item
  Causal models add richness: exploration can target \emph{structural}
  uncertainty, not just parameter uncertainty
\end{itemize}
\end{frame}

\begin{frame}{Causal Multi-Armed Bandits}
\phantomsection\label{causal-multi-armed-bandits}
\begin{itemize}
\tightlist
\item
  Classic multi-armed bandit: \(K\) arms with unknown reward
  distributions

  \begin{itemize}
  \tightlist
  \item
    At each step, pull one arm, observe reward, update beliefs
  \item
    Objective: minimize \textbf{\textcolor{red}{regret}} relative to the
    best arm
  \end{itemize}
\item
  \textbf{\textcolor{darkgreen}{Causal bandit}}: arms are
  \emph{interventions} \(\operatorname{do}(a)\) on a known (or partly
  known) causal graph

  \begin{itemize}
  \tightlist
  \item
    Observing one arm provides information about \emph{other} arms via
    the graph
  \item
    Much lower regret than a standard bandit in many settings
  \end{itemize}
\item
  \textbf{\textcolor{blue}{Example}}: advertising

  \begin{itemize}
  \tightlist
  \item
    Arms: combinations of channel, audience, creative
  \item
    Causal structure: audience and creative \emph{interact}, channel has
    additive effect
  \item
    Exploiting this structure cuts samples needed to find the best
    combination
  \end{itemize}
\end{itemize}
\end{frame}

\begin{frame}{Thompson Sampling}
\phantomsection\label{thompson-sampling}
\begin{itemize}
\tightlist
\item
  \textbf{\textcolor{red}{Thompson sampling}}: a Bayesian exploration
  strategy

  \begin{itemize}
  \tightlist
  \item
    Maintain a posterior over each arm's mean reward
  \item
    At each step, sample a parameter from each posterior and pull the
    arm with the highest sampled mean
  \end{itemize}
\item
  \textbf{\textcolor{darkgreen}{Algorithm}}:

  \begin{enumerate}
  \tightlist
  \item
    Start with prior \(\Pr(\theta_a)\) for each arm \(a\)
  \item
    Sample \(\tilde\theta_a \sim \Pr(\theta_a | \mathcal{D})\) for each
    \(a\)
  \item
    Pull \(a^* = \arg\max_a \tilde\theta_a\)
  \item
    Observe reward, update posterior, repeat
  \end{enumerate}
\item
  \textbf{\textcolor{blue}{Properties}}:

  \begin{itemize}
  \tightlist
  \item
    Naturally balances exploration and exploitation
  \item
    Near-optimal regret bounds \(O(\sqrt{T K \log T})\)
  \item
    Easy to adapt to causal structure by sampling from posteriors over
    causal effects
  \end{itemize}
\end{itemize}
\end{frame}

\begin{frame}{}
\phantomsection\label{section-15}
\begin{itemize}
\tightlist
\item
  Why Prediction Is Not Enough
\item
  Foundations of Decision Theory
\item
  Decision Support with Causal Models
\item
  Sequential Decision-Making and Active Learning

  \begin{itemize}
  \tightlist
  \item
    Value of Information
  \item
    Exploration vs.~Exploitation
  \item
    \emph{\textbf{\textcolor{red}{Bayesian Optimization for Experimentation}}}
  \end{itemize}
\item
  Uncertainty in Causal Decisions
\end{itemize}
\end{frame}

\begin{frame}{Bayesian Optimization: The Setup}
\phantomsection\label{bayesian-optimization-the-setup}
\begin{itemize}
\tightlist
\item
  \textbf{\textcolor{red}{Bayesian optimization (BO)}} targets problems
  where:

  \begin{itemize}
  \tightlist
  \item
    The objective \(f(x)\) is \textbf{\textcolor{orange}{expensive}} to
    evaluate (one experiment per evaluation)
  \item
    No gradients available
  \item
    \(x\) lies in a low- to moderate-dimensional space
  \end{itemize}
\item
  \textbf{\textcolor{darkgreen}{Recipe}}:

  \begin{enumerate}
  \tightlist
  \item
    Place a prior over \(f\), typically a
    \textbf{\textcolor{teal}{Gaussian process}}
  \item
    At each round, fit the GP posterior on observed \((x_i, f(x_i))\)
  \item
    Select \(x_{next}\) maximizing an
    \textbf{\textcolor{cyan}{acquisition function}}
  \item
    Query \(f(x_{next})\), update, repeat
  \end{enumerate}
\item
  \textbf{\textcolor{blue}{Example use cases}}:

  \begin{itemize}
  \tightlist
  \item
    Hyperparameter tuning for a deep model
  \item
    Tuning web page layout
  \item
    Picking reagent concentrations in a chemistry lab
  \end{itemize}
\end{itemize}
\end{frame}

\begin{frame}{Acquisition Functions}
\phantomsection\label{acquisition-functions}
\begin{itemize}
\tightlist
\item
  The acquisition function trades off exploration and exploitation:

  \begin{itemize}
  \tightlist
  \item
    \textbf{\textcolor{red}{Expected improvement (EI)}}:
    \(\alpha_{EI}(x) = \EE[\max(f(x) - f^*, 0)]\)
  \item
    \textbf{\textcolor{orange}{Upper confidence bound (UCB)}}:
    \(\alpha_{UCB}(x) = \mu(x) + \kappa \sigma(x)\)
  \item
    \textbf{\textcolor{darkgreen}{Thompson sampling}}: sample
    \(\tilde f \sim \Pr(f | \mathcal{D})\), maximize \(\tilde f\)
  \item
    \textbf{\textcolor{teal}{Entropy search}}: pick \(x\) that most
    reduces entropy of \(\arg\max f\)
  \end{itemize}
\item
  \textbf{\textcolor{cyan}{Example}}: tuning a learning rate

  \begin{itemize}
  \tightlist
  \item
    Posterior mean peaks at \(10^{-3}\), variance is largest at
    \(10^{-5}\)
  \item
    EI may query \(10^{-5}\) if improvement potential \textgreater{}
    current best
  \item
    UCB with high \(\kappa\) behaves similarly
  \end{itemize}
\item
  \textbf{\textcolor{blue}{Key takeaway}}: the acquisition function
  \emph{is} the policy, derived from the posterior and a utility for
  information
\end{itemize}
\end{frame}

\begin{frame}{BO in Causal Experimentation}
\phantomsection\label{bo-in-causal-experimentation}
\begin{itemize}
\tightlist
\item
  Causal BO combines BO with causal graphs:

  \begin{itemize}
  \tightlist
  \item
    Decision variable \(x\) is an \textbf{\textcolor{red}{intervention}}
    \(\operatorname{do}(X = x)\)
  \item
    Surrogate model is a
    \textbf{\textcolor{orange}{structural causal model}}, not just a GP
    on \(y\)
  \item
    Leverages the graph to share information across interventions
  \end{itemize}
\item
  \textbf{\textcolor{darkgreen}{Example}}: drug dose optimization

  \begin{itemize}
  \tightlist
  \item
    Standard BO models \(y = f(\text{dose})\) directly
  \item
    Causal BO models intermediate biomarkers that mediate dose \(\to\)
    outcome
  \item
    Fewer trials needed when the causal mechanism is partly known
  \end{itemize}
\item
  \textbf{\textcolor{teal}{Pros}}: sample efficient, principled
  exploration, uncertainty-aware
\item
  \textbf{\textcolor{cyan}{Cons}}: requires accurate causal structure;
  misspecification can hurt more than help
\end{itemize}
\end{frame}

\begin{frame}{}
\phantomsection\label{section-16}
\begin{itemize}
\tightlist
\item
  Why Prediction Is Not Enough
\item
  Foundations of Decision Theory
\item
  Decision Support with Causal Models
\item
  Sequential Decision-Making and Active Learning
\item
  \emph{\textbf{\textcolor{red}{Uncertainty in Causal Decisions}}}

  \begin{itemize}
  \tightlist
  \item
    Aleatoric Uncertainty
  \item
    Epistemic Uncertainty
  \item
    Communicating Uncertainty
  \end{itemize}
\end{itemize}
\end{frame}

\begin{frame}{}
\phantomsection\label{section-17}
\begin{itemize}
\tightlist
\item
  Why Prediction Is Not Enough
\item
  Foundations of Decision Theory
\item
  Decision Support with Causal Models
\item
  Sequential Decision-Making and Active Learning
\item
  Uncertainty in Causal Decisions

  \begin{itemize}
  \tightlist
  \item
    \emph{\textbf{\textcolor{red}{Aleatoric Uncertainty}}}
  \item
    Epistemic Uncertainty
  \item
    Communicating Uncertainty
  \end{itemize}
\end{itemize}
\end{frame}

\begin{frame}{Aleatoric Uncertainty: Irreducible Randomness}
\phantomsection\label{aleatoric-uncertainty-irreducible-randomness}
\begin{itemize}
\tightlist
\item
  \textbf{\textcolor{red}{Aleatoric uncertainty}} comes from the
  \emph{inherent} randomness in the outcome

  \begin{itemize}
  \tightlist
  \item
    Also called \emph{statistical} or \emph{risk} uncertainty
  \item
    Cannot be reduced by collecting more data
  \item
    Captured by \(\Pr(Y | \operatorname{do}(a), \theta)\) given known
    \(\theta\)
  \end{itemize}
\item
  \textbf{\textcolor{darkgreen}{Example}}: flipping a known-fair coin

  \begin{itemize}
  \tightlist
  \item
    \(\Pr(\text{heads}) = 0.5\) exactly
  \item
    More flips don't sharpen the single-flip outcome
  \end{itemize}
\item
  \textbf{\textcolor{blue}{Example}}: medical treatment with a known
  response distribution

  \begin{itemize}
  \tightlist
  \item
    Even if we knew the drug's exact mechanism, individual patient
    outcomes would vary due to biological stochasticity
  \end{itemize}
\item
  \textbf{\textcolor{violet}{Decision implication}}: aleatoric
  uncertainty sets the \emph{floor} on risk --- no amount of learning
  removes it
\end{itemize}
\end{frame}

\begin{frame}{Representing Aleatoric Uncertainty}
\phantomsection\label{representing-aleatoric-uncertainty}
\begin{itemize}
\tightlist
\item
  In a structural causal model: \[
  Y = f(\text{Pa}(Y), U_Y), \quad U_Y \sim \Pr(U_Y)
  \]

  \begin{itemize}
  \tightlist
  \item
    \(U_Y\) is the exogenous noise capturing aleatoric variability
  \end{itemize}
\item
  Common distributional assumptions:

  \begin{itemize}
  \tightlist
  \item
    Gaussian: \(U_Y \sim \mathcal{N}(0, \sigma^2)\) for continuous \(Y\)
  \item
    Bernoulli / categorical for discrete outcomes
  \item
    Heavy-tailed (Student-\(t\), Laplace) for robust modeling
  \end{itemize}
\item
  \textbf{\textcolor{red}{Q}}: How do we tell the agent apart from the
  noise when we only see \(Y\)?

  \begin{itemize}
  \tightlist
  \item
    We can't, unless \(U_Y\) has a tractable parametric form and we
    exploit the causal structure
  \end{itemize}
\item
  \textbf{\textcolor{blue}{Example}}: a recommender's click-through rate
  varies user-to-user; even with perfect knowledge of
  \(\Pr(\text{click} | \text{user, item})\), specific clicks are random
\end{itemize}
\end{frame}

\begin{frame}{}
\phantomsection\label{section-18}
\begin{itemize}
\tightlist
\item
  Why Prediction Is Not Enough
\item
  Foundations of Decision Theory
\item
  Decision Support with Causal Models
\item
  Sequential Decision-Making and Active Learning
\item
  Uncertainty in Causal Decisions

  \begin{itemize}
  \tightlist
  \item
    Aleatoric Uncertainty
  \item
    \emph{\textbf{\textcolor{red}{Epistemic Uncertainty}}}
  \item
    Communicating Uncertainty
  \end{itemize}
\end{itemize}
\end{frame}

\begin{frame}{Epistemic Uncertainty: Model Misspecification}
\phantomsection\label{epistemic-uncertainty-model-misspecification}
\begin{itemize}
\tightlist
\item
  \textbf{\textcolor{red}{Epistemic uncertainty}} stems from \emph{lack
  of knowledge}

  \begin{itemize}
  \tightlist
  \item
    Also called \emph{systematic} or \emph{knowledge} uncertainty
  \item
    Reducible in principle by collecting more/better data or improving
    the model
  \item
    Captured by the posterior \(\Pr(\theta | \mathcal{D})\)
  \end{itemize}
\item
  \textbf{\textcolor{darkgreen}{Sources}}:

  \begin{itemize}
  \tightlist
  \item
    Limited or biased data
  \item
    Wrong functional form
  \item
    Missing variables or wrong causal graph
  \item
    Confounders not adjusted for
  \end{itemize}
\item
  \textbf{\textcolor{blue}{Example}}: a small clinical trial with
  \(n = 40\) patients

  \begin{itemize}
  \tightlist
  \item
    Wide posterior on treatment effect \(\theta\)
  \item
    Gathering more data shrinks this posterior
  \end{itemize}
\end{itemize}
\end{frame}

\begin{frame}{Aleatoric vs.~Epistemic: Side-by-Side}
\phantomsection\label{aleatoric-vs.-epistemic-side-by-side}
\begin{columns}[T]
\begin{column}{0.5\textwidth}
\textbf{\textcolor{red}{Aleatoric}}

\begin{itemize}
\item
  Source: inherent randomness
\item
  Reducible: \textbf{\textcolor{orange}{No}}
\item
  Represented: noise \(U\)
\item
  Grows with: complexity of \(Y\)
\item
  Managed by: risk preferences
\item
  Example: coin flip, weather
\item
  \textbf{\textcolor{darkgreen}{Implication}}: influences utility
  function shape
\end{itemize}
\end{column}

\begin{column}{0.5\textwidth}
\textbf{\textcolor{teal}{Epistemic}}

\begin{itemize}
\item
  Source: model \& data limits
\item
  Reducible: \textbf{\textcolor{cyan}{Yes}}, via evidence
\item
  Represented: posterior on \(\theta\)
\item
  Shrinks with: data, experiments
\item
  Managed by: learning, VOI
\item
  Example: effect of a new drug
\item
  \textbf{\textcolor{blue}{Implication}}: motivates sequential
  decision-making
\end{itemize}
\end{column}
\end{columns}

\begin{itemize}
\tightlist
\item
  \textbf{\textcolor{violet}{Key takeaway}}: good decision systems
  \emph{separate} the two --- combining them collapses information that
  matters for action
\end{itemize}
\end{frame}

\begin{frame}{Distinguishing in Practice}
\phantomsection\label{distinguishing-in-practice}
\begin{itemize}
\item
  A Bayesian predictive distribution decomposes as: \[
  \Pr(Y | \operatorname{do}(a), \mathcal{D}) =
  \int \Pr(Y | \operatorname{do}(a), \theta) \, \Pr(\theta | \mathcal{D}) \, d\theta
  \]
\item
  \textbf{\textcolor{red}{Variance decomposition}}: \[
  \VV[Y | \operatorname{do}(a), \mathcal{D}] =
  \underbrace{\EE_\theta[\VV[Y | a, \theta]]}_{\text{aleatoric}} +
  \underbrace{\VV_\theta[\EE[Y | a, \theta]]}_{\text{epistemic}}
  \]
\item
  \textbf{\textcolor{darkgreen}{Example}}: deep ensemble for sales
  forecasting

  \begin{itemize}
  \tightlist
  \item
    Within-model variance \(\approx\) aleatoric
  \item
    Between-model variance \(\approx\) epistemic
  \item
    Report both to the decision-maker
  \end{itemize}
\item
  \textbf{\textcolor{blue}{Warning}}: most off-the-shelf ML models
  report only a point prediction --- neither kind of uncertainty is
  visible
\end{itemize}
\end{frame}

\begin{frame}{}
\phantomsection\label{section-19}
\begin{itemize}
\tightlist
\item
  Why Prediction Is Not Enough
\item
  Foundations of Decision Theory
\item
  Decision Support with Causal Models
\item
  Sequential Decision-Making and Active Learning
\item
  Uncertainty in Causal Decisions

  \begin{itemize}
  \tightlist
  \item
    Aleatoric Uncertainty
  \item
    Epistemic Uncertainty
  \item
    \emph{\textbf{\textcolor{red}{Communicating Uncertainty}}}
  \end{itemize}
\end{itemize}
\end{frame}

\begin{frame}{Communicating Uncertainty to Stakeholders}
\phantomsection\label{communicating-uncertainty-to-stakeholders}
\begin{itemize}
\item
  Stakeholders often prefer a single number; giving one is a disservice
\item
  \textbf{\textcolor{red}{Best practices}}:

  \begin{itemize}
  \tightlist
  \item
    Report \textbf{\textcolor{orange}{intervals}} (90\% posterior
    credible intervals) alongside point estimates
  \item
    Distinguish aleatoric from epistemic in words and visuals
  \item
    Frame probabilities in natural frequencies: \emph{``about 2 in
    100''}
  \item
    Use \textbf{\textcolor{darkgreen}{visual aids}}: fan charts, violin
    plots, calibration plots
  \item
    Link uncertainty to \textbf{\textcolor{teal}{actions}}: \emph{``if
    we get 200 more samples, the interval shrinks from \$5M to \$1M''}
  \end{itemize}
\item
  \textbf{\textcolor{cyan}{Anti-patterns}}:

  \begin{itemize}
  \tightlist
  \item
    Reporting only point predictions
  \item
    Combining aleatoric and epistemic into one opaque bar
  \item
    Hiding priors in appendices
  \end{itemize}
\end{itemize}
\end{frame}

\begin{frame}{Communicating Uncertainty: A Worked Example}
\phantomsection\label{communicating-uncertainty-a-worked-example}
\begin{itemize}
\item
  \textbf{\textcolor{red}{Setting}}: pricing team evaluating a new fee
  structure
\item
  \textbf{\textcolor{darkgreen}{Bad communication}}:

  \begin{itemize}
  \tightlist
  \item
    \emph{``The model predicts \$2.3M in additional revenue.''}
  \end{itemize}
\item
  \textbf{\textcolor{blue}{Better communication}}:

  \begin{itemize}
  \tightlist
  \item
    \emph{``Our best estimate is \$2.3M. 90\% of the uncertainty sits
    between \$0.8M and \$3.7M.}
  \item
    \emph{Of this, about \$1.5M of the width is intrinsic customer
    variability (aleatoric) --- a pilot won't narrow that much.}
  \item
    \emph{The remaining \$1.4M is parameter uncertainty (epistemic) and
    would shrink by roughly 60\% if we ran the proposed 4-week pilot.''}
  \end{itemize}
\item
  \textbf{\textcolor{violet}{Decision impact}}: the pilot discussion now
  has a clear expected value, not a gut feeling
\end{itemize}
\end{frame}

\begin{frame}{Summary: Decision-Making with Causal Models}
\phantomsection\label{summary-decision-making-with-causal-models}
\begin{itemize}
\item
  \textbf{\textcolor{red}{Prediction alone is insufficient}} for
  choosing actions --- Simpson's paradox, confounding, and policy
  reversal all require causal reasoning
\item
  \textbf{\textcolor{orange}{Expected utility}} under interventional
  distributions is the right target \[
  a^* = \arg\max_a \EE_{Y \sim \Pr(Y | \operatorname{do}(a))}[U(Y)]
  \]
\item
  \textbf{\textcolor{darkgreen}{Influence diagrams}} extend causal DAGs
  with decisions and utilities, giving a graphical, solvable decision
  problem
\item
  \textbf{\textcolor{teal}{Sequential decisions}} need exploration: VOI
  tells us when to gather information; Thompson sampling and Bayesian
  optimization tell us \emph{how}
\item
  \textbf{\textcolor{cyan}{Uncertainty}} decomposes into aleatoric
  (irreducible) and epistemic (reducible) --- both must be modeled and
  communicated
\item
  \textbf{\textcolor{blue}{Key takeaway}}: causal reasoning, decision
  theory, and Bayesian inference are three sides of the same coin ---
  and together they turn predictions into \emph{robust} actions
\end{itemize}
\end{frame}

\end{document}
